\chapter{Miscelleanous}
\begin{itemize}
\item Please try to reuse already existing code to avoid code duplication. For example, your new crypto function uses a \enquote{Polybios" cipher} internally, it is a good idea to reuse the Polybios code for creating the grid and use the encryption/decryption. E.g. have a look at the ADFGVX or Bifid Ciphers. Many functions are included in the code. There's a great chance that something similar to your new function has already been done.
\item Other basic crypto functions are \enquote{Transposition} \enquote{logic/tools/crypto\_end\_encodings/transposition.dart}; moves characters or strings from rows to columns and vice versa) and \enquote{Substitution} \enquote{logic/tools/crypto\_and\_encodings/substitution.dart}; very common operation in many ciphers; replaces characters or strings with new characters or strings)
\item When it comes to coordinate calculations, the base functions \enquote{logic/tools/coords/projection/projection()} for waypoint projection and \enquote{logic/tools/coords/distance\_and\_bearing/distanceBearing()} for calculation of distances and bearings between two coordinates should be useful.
\item Maybe you discover new possibilities for centralizing some functions or values, please do not hesitate to encapsulate them to a more common place. Cleanup and refactoring is an important and necessary step in coding. For example, you could use one of the Utils files for such centralization if they fit.
\item Please avoid using Frontend/UI classes or types in logic or test code. This avoids many nested internal Flutter UI dependencies where there are not necessary. These types usually have the prefix <tt>GCW</tt> and they are located in the <tt>widgets</tt> path. Best way to avoid this, is to write the logic/test code before creating the UI. If you need, for example, the output of a Switch widget (usually of type <tt>GCWSwitchPosition</tt>), do not hesitate to map it to a more common enum type, that can be used easily within the logic or test code.
\item Please always use spaces around comparators and operators: \lstinline|var i = a + b|
\item The Dart coding style uses two spaces as default indent. Please use this, althought it may looks nasty... ;)
\item The localisation keywords currently follow their own rules: E.g. let's take \enquote{coords\_formatconverter\_mgrs\_easting}: This keywords contains four parts, separated by an underscore. Those can be seen as an hierachy: Topmost the coords section, after that the tool \enquote{Format Converter} here without any separator and without camel case!, afterwards the sub tool \enquote{MGRS}, which has an \enquote{Easting}. So, please try to keep such an order when creating a tool. I believe, in more complex cases a separator for parts could make sense. Then please use a dot: \enquote{coords\_format.converter\_mgrs\_easting}
\item Please name all methods and variables in English. All comments have to be in English, as well.
\end{itemize}
